\documentclass[tikz,border=10pt]{standalone}
\usepackage{fontspec}
\usepackage{xeCJK}
\setmainfont{Times New Roman}
\setCJKmainfont{Hiragino Sans GB}
\setCJKsansfont{Microsoft YaHei}
\setCJKmonofont{STSong}
\usepackage{tikz}
\usetikzlibrary{arrows.meta,matrix,positioning}

\tikzset{
  axis/.style={-{Stealth[length=3mm]}, line width=1.1pt, draw=black!75},
  rl/.style={draw=red!75, line width=1.4pt},
  pole/.style={draw=red!75, line width=1.2pt},
  zero/.style={draw=red!75, line width=1.2pt},
  cell/.style={draw=black!70, line width=1.0pt, minimum width=5.1cm, minimum height=1.05cm, align=center, text width=4.6cm},
  head/.style={cell, draw=blue!70!black, line width=1.2pt},
  left/.style={cell, draw=green!50!black, line width=1.2pt}
}

\begin{document}
\begin{tikzpicture}
  \begin{scope}[shift={(0,0)}]
    \draw[axis] (-2.8,0) -- (2.8,0) node[right] {$\sigma$};
    \draw[axis] (0,-2.6) -- (0,2.9) node[above] {$j\omega$};
    \draw[rl] (-2,0) -- (-1,0);
    \draw[rl] (-1,0) -- (-1,2.2);
    \draw[rl] (-1,0) -- (-1,-2.2);
    \draw[pole] (-2.15,-0.15) -- (-1.85,0.15);
    \draw[pole] (-2.15,0.15) -- (-1.85,-0.15);
    \draw[pole] (-0.15,-0.15) -- (0.15,0.15);
    \draw[pole] (-0.15,0.15) -- (0.15,-0.15);
    \fill[green!60!black] (-1,0) circle (0.07);
    \node[below] at (-2,0) {$-2$};
    \node[below] at (0,0) {$0$};
    \node[below left] at (-1,0) {$-1$};
    \node[blue!70!black] at (-0.2,2.4) {$K^\ast \uparrow$};
  \end{scope}

  \begin{scope}[xshift=7.0cm]
    \matrix[matrix of nodes, nodes in empty cells, row sep=-\pgflinewidth, column sep=-\pgflinewidth] {
      \node[head] {视角}; & \node[head] {工程含义}; \\
      \node[left] {动态行为}; & \node[cell] {先是实极点，后转为共轭复根；系统由过阻尼走向振荡}; \\
      \node[left] {稳定性}; & \node[cell] {$\Re\{\lambda_{1,2}\}<0$ 在全部图示分支上都成立，因此闭环始终绝对稳定}; \\
      \node[left] {稳态误差}; & \node[cell] {对斜坡输入有 $e_{ss}=2A/K^\ast$，因此 $K^\ast$ 越大，跟踪误差越小}; \\
    };
  \end{scope}
\end{tikzpicture}
\end{document}
